\subsection*{整数与元组的使用规则}

一个控制值用整数表示，一组有顺序的值用 \co@ituple@ 表示。形状 \co@[7, 14]@ 描述范围，元组 \co@{1, 2}@ 则可以表示位置或分块因子。两者的秩一致时，形状运算会把每个因子应用到对应维度。

整数和普通元组都必须初始化，之后不能重新赋值。这项限制针对用于控制和形状计算的名字。缓冲区元素则不同：加法示例会写入 \co@output.at(i, j)@，矩阵乘法也会更新累加器。区分这两种角色，就不会把不可变索引误当成可变累加器。
